%% kapa: exam=1 round=30 year=2019 period=1 form=A subject=real_estate subject_ko=부동산학원론 seq=15 number=95
%% source: Q-Net 감정평가사 기출문제 (cst00309) artlSeq=5205205 file=extracted/1008/5205205/2019년 제30회 감정평가사 1차시험 문제지/2019년 제30회 감정평가사 1차시험 문제지/2019년도 제30회 감정평가사 1차 1교시 A형.hwp
%% conversion: kordoc 4.24.51 -> tools/gen1.py (auto); review=proofread-sonnet; flags=none
\begin{question}{95}
사업주가 특수목적회사인 프로젝트 회사를 설립하여 특정 프로젝트 수행에 필요한 자금을 금융기관으로부터 대출받는 방식의 프로젝트 금융을 활용하는 경우에 관한 설명으로 옳지 않은 것은? (단, 프로젝트 회사를 위한 별도의 보증이나 담보제공 등은 없음)
\begin{choices}
  \item 대규모 자금이 소요되고 공사기간이 장기인 사업에 적합한 자금조달수단이다.
  \item 프로젝트 금융에 의한 채무는 사업주와 독립적이므로 부채상환의무가 사업주에게 전가되지 않는다.
  \item 사업주가 이미 대출한도를 넘어섰거나 대출제약요인이 있는 경우에도 가능하다.
  \item 해당 프로젝트가 부실화되더라도 대출기관의 채권회수에는 영향이 없다.
  \item 프로젝트 회사는 법률적, 경제적으로 완전히 독립적인 회사이지만 이해당사자간의 이견이 있을 경우에는 사업지연을 초래할 수 있다.
\end{choices}
\end{question}
\answer{4}
